% !TeX root = ../papers/nc-p3.tex
\section{有限窓の反対階数九における片側対称性}
\label{sec:rank-nine-finite-symmetry}
\begin{theorem}\label{thm:rank-nine-finite-symmetry}
長さ1の局所点を含む有限Frobenius点グラフで、反対階数が9とする。
中支持3の点による認証がないとき、窓は少なくとも一方の対向位置の交換で不変な表示を持つ。
両対称なら既出の有限AS判定を使える。片側対称表示の対称二乗の核は一つの二次形式で、
その階数が2以下なら全次数で正則である。従ってこの範囲の未判定窓では、核の階数は3か4に限られる。
\end{theorem}

\begin{lemma}\label{lem:irreducible-four-plane-brackets}
三行跡零行列の既約な四次元部分空間 $D$ の交換子の線型包は三次元以上である。
\end{lemma}
\begin{proof}
$D$ に相異なる三固有値の行列があることをまず示す。
反対なら各元 $X$ の三次特性式の判別式が零で、$\operatorname{tr}(X^2)^3$ と
$\operatorname{tr}(X^3)^2$ は恒等的に一定の非零倍率で一致する。
ともに恒等零か、$D$ の座標環の一意分解で共通の一次式の二乗・三乗の非零倍である。
従って次元3以上の冪零部分空間 $N$ がある。三行の冪零部分空間で次元3以上なら
厳密上三角全体に同時共役であることをここで確認する。
巡回冪零元がないと全元が階数1以下で、共通像直線または共通核平面および跡零条件から次元は2以下である。
巡回元を $E_{12}+E_{23}$ とすれば、これと $Y\in N$ の全線型結合の平方・三乗の跡が零なので
$Y_{31}=0,Y_{32}=-Y_{21}$、$Y_{21}(Y_{11}-Y_{33})=0$ となる。
$Y_{21}$ が恒等零でなければ $Y_{11}=Y_{33}=a,\ Y_{21}=c$ と書き
$3a^2+c(Y_{12}-Y_{23})=0$。よって $a=kc,\ Y_{23}=Y_{12}+3k^2c$ である。
さらに $\det Y=0$ は $c^2(k^3c-Y_{13})=0$ を与える。自由座標は $Y_{12},c$ の二個で矛盾する。
従って $Y_{21}=0$ で全元が上三角冪零である。
$D$ 全体が冪零であることはこれで排除された。残るとき $N$ は一次式の核であり、
平方の跡の偏極より $\operatorname{tr}(D N)=0$。従って $D$ も上三角で既約性に反する。

交換子の包が二次元以下と仮定する。一般の $X\in D$ の $D$ 内の中心化空間は次元2以上となり、
三固有値が異なるので跡零中心化空間の全体である。特に $X^2-\frac13\operatorname{tr}(X^2)I\in D$。
稠密性と偏極で $k I+D$ はJordan積 $(UV+VU)/2$ で閉じる五次元空間になる。
ある三固有値行列を対角化すればその三つの射影冪等元がこの空間に入る。
各冪等元とのJordan積で対角成分と三つの非対角対 $ij,ji$ に分解できる（可換な対角作用の同時重み分解）。
非対角部分の総次元は2である。既約なので、可能な非零行列成分を辺とする三頂点有向グラフは強連結である。
ゆえに相異なる二つの非対角対で各一次元を使い、両辺は双方向でなければならない。
その二元のJordan積は三番目の対に非零成分を持ち、五次元性に反する。
\end{proof}

\begin{proof}[定理の対称性]
印付きの $g=\dim\langle G,F\rangle,\ q=\dim\langle Q,P\rangle$ はともに3以上で $g+q\le7$。
両方3なら前節の証明で $\bar w=0$ を必要としたのは対称・交代の混合型を排除する段だけである。
混合型でも $G$ を交代、$Q$ を対称とすれば $F=G,P=Q$ である。
$a,c$ 交換では星印を含む項は不変、純成分の差は
$[G_a,G_c]_{bd}+[Q_b,Q_d]_{ac}=0$ である。よって少なくとも一方は対称となる。

残りは巡回して $g=4,q=3$ とする。$P_d=Q_{Jd}$ であり、
$G,F$ の三平面は二平面で交わる。このため $D=\langle G,F\rangle$ 全体の外積は交叉交換子の入力
$G\wedge F$ で張られる。実際、共通二元とそれぞれの補元による六基底外積を使える。
既約な $D$ に前補題を適用すると、$[G,F]$ の包は三次元以上、一方で $[Q,Q]$ の包と同じ次元なので
ちょうど三次元である。交換子式の $b,d$ 跡から $\sum_b[Q_b,Q_{Jb}]=0$、
独立性から $J$ は対称である。合同変換で $J=I$ と取れば $D$ の全交換子は交代行列となる。
ここで $D$ は全て対称である。実際、交代部分の像次元が1,2なら、前節と同じ中心化空間の議論により
その核に必要な二つ以上の跡零対称行列がない。像次元3なら三元の交代部分を外積基底に取れ、
同じ十五次の検査でその三元は純交代となる。残りの対称元は三外積の全てと交換するので零になってしまう。
よって $G,F$ は対称で、$P=Q$ はこの段では任意の跡零三行列組でよい。
$b,d$ を交換しても星印を持つ項は不変、純成分の差も交換子式で零となる。
これで全ての場合の片側対称性を得た。

対称側の二次核は $10-9=1$ 次元である。これが階数2以下なら $x y+ y x$ 型に書け、
対称性で $w(x,-,y,-)=0$ となる（$x=y$ も許す）。
六次元標的の隣接切断 $\beta,\gamma$ では $\beta(x,-)$ と $\gamma(y,-)$ の像が互いに直交する。
有限ファイバーより両評価は階数3以上なので両方3、従って $x$ は点グラフ上にある。
その三テンソルは中央を $y$ で消せるから中支持3となり既知の全次数認証が適用できる。
\end{proof}
